% Complete proof candidate. Individual check status is in result_audit.json.

\subsection{The Betti measure and torsion equidistribution}

We take $dd^c=(i/\pi)\partial\bar\partial$, so that
$dd^c\log|w|=\delta_0$. For a local complex coordinate $w$,
$dA_w=(i/2)dw\wedge d\bar w$.

\begin{proposition}[ARI C.8.9; Refuted sign and normalization, with the corrected current identity proved]\label{ari:C.8.9}
On a local period chart write $\mu=a\tau+b$, with $a,b$ real,
and put $T=\Im\tau$. The positive Betti measure is
\[
 m_P=db\wedge da
 =\frac{|\mu_w-a\tau_w|^2}{T}\,dA_w.
\]
Thus the form $da\wedge db$ in Lead 15(b) has the wrong sign.
The globally single-valued function $L(t)=\log N(P_t)$ satisfies
on $\mathbb P^1(\C)$
\begin{equation}\label{ari:eq:ca-current}
 dd^cL=\frac1{12}\delta_\infty-m_P,\qquad
               \int_{\mathbb P^1}m_P=\frac1{12}.
\end{equation}
There are no atoms of $m_P$ at $0$, $1$ or $\infty$.
Let $T_0=\sqrt3/2$, $\rho=(1+i\sqrt3)/2$, and
$p=\wp'_\rho((2-\rho)/3)$. The density in the ordinary
$t$ coordinate satisfies
\begin{align}
 \frac{dm_P}{dA_t}
 &=\frac{|t|^{-2/3}}{64T_0|p|^{2/3}}
                      \bigl(1+O(|t|^{1/3})\bigr),
                         &&t\longrightarrow0,\label{ari:eq:ca-density0}\\
 \frac{dm_P}{dA_t}
 &=\frac{|t-1|^{-3/2}}{24\pi^2\sqrt{E_4(i)}}
                      \bigl(1+O(|t-1|^{1/2})\bigr),
                         &&t\longrightarrow1.\label{ari:eq:ca-density1}
\end{align}
\end{proposition}
\begin{proof}
On a chart avoiding a zero of the section, the theta quotient
is holomorphic and nonzero. Its log modulus is pluriharmonic.
Only $-\pi(\Im\mu)^2/T$ contributes to the curvature.
Direct differentiation of this real quotient gives
\[
 \partial_w\partial_{\bar w}\frac{(\Im\mu)^2}{T}
       =\frac{|\mu_w-a\tau_w|^2}{2T}.
\]
Also
\[
 a_w=\frac{\mu_w-a\tau_w}{2iT},\qquad
 b_w=-\bar\tau\frac{\mu_w-a\tau_w}{2iT}.
\]
Taking the wedge product proves the formula for $m_P$ and
$dd^cL=-m_P$ on the chart. The transformation of $(a,b)$
under an integral change of period basis has determinant one,
so the measure patches. Translation of $\mu$ by a locally
constant integral period changes neither the formula nor
the measure. The formula remains valid on a root chart
at an elliptic modulus.

At $t=0$ the expansion proved in C.3 is
\[
 L(t)=\frac14\log3+\frac7{192}\Re t
       -B|t|^{4/3}+O(|t|^{5/3}),\qquad
 B=\frac{9\pi}{512T_0|p|^{2/3}}.
\]
The expansion is a convergent real-analytic expansion on a
root chart, so it can be differentiated termwise with the
indicated orders. The Euclidean identity
$\Delta_{\R^2}|t|^\alpha=\alpha^2|t|^{\alpha-2}$ gives
\eqref{ari:eq:ca-density0}.

At $t=1$ put $h=\mu-(1-\tau)/2$. The checked K3 expansion
and the modular $j$ expansion give
\[
 t-1=A_1h^4(1+O(h^2)),\qquad
 A_1=\frac94\pi^4E_4(i)>0,\qquad
 L=\log\sqrt2-\frac\pi2|h|^2+O(|h|^4).
\]
Therefore
$L=\log\sqrt2-\pi|t-1|^{1/2}/(2\sqrt{A_1})+O(|t-1|)$.
Termwise differentiation on the root chart proves
\eqref{ari:eq:ca-density1}. Both densities are locally integrable,
and the continuous potentials at $0$ and $1$ have no atoms.

For infinity take $u=1/t$, and set $\ell=\log(5+2\sqrt6)$.
The cusp branch is
$\mu=1/2-i\ell/(2\pi)+O(q)$, $q=e^{2\pi i\tau}$.
The theta and eta products, with the nonzero factor
$|2\sin\pi\mu|=2\sqrt3+O(q)$, give
\begin{equation}\label{ari:eq:ca-cuspexp}
 L=\frac1{12}\log|u|+\log(2\sqrt3)-\frac1{12}\log1728
          -\frac{\ell^2}{2\log(1728/|u|)}
          +O\left(\frac{|u|}{|\log|u||}+|u|\right).
\end{equation}
Here $q=u/1728+O(u^2)$; the estimate follows directly by
bounding the convergent products. The nonlogarithmic
correction is continuous at $u=0$ and has no point mass.
Equivalently its Betti density is
\[
 \frac{dm_P}{dA_u}
  =\frac{\ell^2}{2\pi|u|^2\log(1728/|u|)^3}
       \bigl(1+o(1)\bigr),
\]
which is integrable. The single logarithmic term in
\eqref{ari:eq:ca-cuspexp} contributes $\delta_\infty/12$.
This proves the global current identity; pairing it with
the constant function one gives the total mass.
\end{proof}

\begin{proposition}[ARI C.8.10; Proved: the correctly normalized equidistribution]\label{ari:C.8.10}
Let $D_n$ be the divisor, with intersection multiplicity, of
parameters $t\in\mathbb P^1\setminus\{0,1,\infty\}$ for which
$nP_t=O$. Then
\[
 \deg D_n=\frac{n^2}{12}+O(1),\qquad
 \frac1{n^2}[D_n]\ \longrightarrow\ m_P,\qquad
 \frac1{\deg D_n}[D_n]\ \longrightarrow\ 12m_P
\]
weakly as $n\to\infty$ through indices for which $D_n\ne0$.
For any nonrepeating sequence of individual algebraic torsion
parameters, their Galois-orbit probability measures likewise
converge to $12m_P$. Thus the probability-limit assertion
in Lead 15(b) requires the factor $12$.
\end{proposition}
\begin{proof}
We apply \citet[Theorem 1.1 and Corollary 1.2]{demarco2017}.
Their hypotheses are a number field, a smooth projective
base curve, an elliptic surface and a rational section of
nonzero generic canonical height. Here they are $\Q$,
$\mathbb P^1$, the displayed rational elliptic surface,
and $P$, whose generic height for $(O)$ is $1/12$ by C.6.4.
At a torsion specialization the canonical height is zero.
Corollary 1.2 consequently gives convergence of every
nonrepeating sequence of Galois-orbit probability measures.

The limiting measure is the curvature of the height metric
divided by the degree of its line bundle. In our normalization,
C.8.9 computes the curvature as $m_P$, and C.6.4 computes
that degree as $1/12$. Hence the probability measure is
$12m_P$. This division by the degree is essential: a measure
of mass $1/12$ cannot be a weak limit of probability measures.
The source's informal description of the pullback current
is interpreted with this normalization.

For completeness, the passage from individual orbits to
the full division divisor does not require uniformity of
the equidistribution theorem in $n$. The surface height
formula, with our half-self-pairing convention, reads
\[
 \widehat h_O(Q)=\chi(\mathcal O_{\mathcal E})+(Q\cdot O)
                -\frac12\sum_v\operatorname{contr}_v(Q).
\]
Here $\chi=1$, and the finitely many fibre corrections are
bounded as $Q=nP$ varies, because they depend only on its
component. Since $\widehat h_O(nP)=n^2/12$, it follows
that $(nP\cdot O)=n^2/12+O(1)$. Removing the singular fibres
changes this by $O(1)$: the local intersection at each
such fibre is bounded as $n$ varies in characteristic zero.
One can see this in the formal group after a fixed
multiple enters it; multiplication by a nonzero integer
is a local analytic isomorphism over the residue field
$\C$. Thus $\deg D_n=n^2/12+O(1)$.

At any fixed smooth torsion parameter $t_0$, of exact order
$d$, the multiplicity in $D_n$ is zero unless $d\mid n$,
and is constant when $d\mid n$. Indeed the formal parameter
of $nP=[n/d](dP)$ has first term $(n/d)$ times that of $dP$,
with nonzero constant coefficient. Fix a continuous test
function and $\varepsilon>0$. Corollary 1.2 implies that
all but finitely many distinct Galois orbits of torsion
parameters have average within $\varepsilon$ of the limiting
integral; otherwise a countersequence would violate it.
The total weighted contribution of the finitely many
exceptional orbits is bounded independently of $n$, by
the multiplicity observation. Divide by $\deg D_n$ and
let $n\to\infty$. This proves the probability convergence
for $D_n$, and the degree asymptotic proves the unnormalized
version.
\end{proof}

\begin{proposition}[ARI C.8.11; Proved: the two first-derivative identities]\label{ari:C.8.11}
On an ordinary modulus chart put
\[
 a=\frac{\Im\mu}{T},\qquad
 W=\frac{\vartheta_1'(\mu,\tau)}{\vartheta_1(\mu,\tau)}
                         +2\pi ia,\qquad
 R=\frac{d\mu}{d\tau}-a.
\]
Then
\[
 L_\tau=\frac{W}{8\pi i}(W+4\pi iR),\qquad
 \det\frac{d}{d\tau}
 \begin{pmatrix}\tau&\mu\\\mu&\beta^{\rm part}/6\end{pmatrix}
       =\frac{(W+2\pi iR)^2}{4\pi^2}.
\]
Moreover
\[
 R=\frac{g_2}{6\pi i\wp'(\mu)}-\frac{W}{2\pi i},
 \qquad
 L_\tau=\frac{iW(3W\wp'(\mu)-2g_2)}
                         {24\pi\wp'(\mu)}.
\]
The critical-point factors in K4 follow with their stated
normalization after multiplying by $6/\pi$.
\end{proposition}
\begin{proof}
Put $f=\vartheta_1'/\vartheta_1$ and
$E=\pi^2E_2/3$. The theta heat equation and
$\eta'/\eta=\pi iE_2/12$ give
\[
 \partial_\tau\log(\vartheta_1/\eta)
       =\frac{f_z+f^2+E}{4\pi i}.
\]
At $\wp(\mu)=0$, the identity $f_z=-\wp-E$ reduces
this to $f^2/(4\pi i)$. Thus, along the section,
\[
 \frac{d}{d\tau}\log(\vartheta_1(\mu,\tau)/\eta)
       =f\mu'+\frac{f^2}{4\pi i}.
\]
Differentiate the Gaussian $-\pi(\Im\mu)^2/T$ using
Wirtinger derivatives, retaining the factor $1/2$ in
the derivative of the real part of a holomorphic logarithm.
Substitution of $f=W-2\pi ia$ and $\mu'=R+a$ gives the
first formula. Since
$\beta^{\rm part}/6=(i/\pi)\log(\vartheta_1/(2\sqrt3\eta))$,
its derivative is $if\mu'/\pi+f^2/(4\pi^2)$.
Subtracting $\mu'^2$ gives
$(f+2\pi i\mu')^2/(4\pi^2)$, the second formula.

Finally differentiate $\wp_\tau(\mu(\tau))=0$ by the theta
heat equation. After substituting $\wp(\mu)=0$,
$\wp''(\mu)=-g_2/2$, and $f=W-2\pi ia$, the result is
$6\pi i\wp'R=g_2-3W\wp'$. This is the formula for $R$.
Inserting it into the first identity gives the last one.
These computations also independently verify the K4 step
used in C.3; its apparent singularity at a half-period is
handled by the root chart rather than division by zero.
\end{proof}

